% =====================================================================
% IEEE Conference Paper
% Topic: Vibration-Based Fault Detection and Diagnosis for Industrial
%        Rotating Machinery Using Spectral Analysis and SVM
% =====================================================================
% !TeX spellcheck = en-US
% !TeX encoding = utf8
% !TeX program = pdflatex
% !BIB program = bibtex

\documentclass[conference,a4paper]{IEEEtran}[2015/08/26]

\usepackage{balance}
\usepackage{iftex}
\usepackage[hyphens]{url}
\makeatletter
\g@addto@macro{\UrlBreaks}{\UrlOrds}
\makeatother

\usepackage{graphicx}
\usepackage{amsmath}
\usepackage{amssymb}
\usepackage{bm}

\usepackage[dvipsnames,table]{xcolor}
\usepackage{booktabs}
\usepackage{multirow}
\usepackage{array}
\usepackage{paralist}

\usepackage{tikz}
\usetikzlibrary{shapes.geometric,arrows.meta,positioning,calc,fit,
                decorations.markings,matrix}
\usepackage{pgfplots}
\pgfplotsset{compat=1.17}
\usepgfplotslibrary{groupplots,colorbrewer}

\usepackage{placeins}

\usepackage[autostyle=true]{csquotes}
\usepackage[square,comma,numbers,sort&compress]{natbib}
\renewcommand{\bibfont}{\normalfont\footnotesize}

\usepackage[group-minimum-digits=4,per-mode=fraction]{siunitx}

\usepackage{hyperref}
\hypersetup{hidelinks,colorlinks=true,allcolors=black,
            pdfstartview=Fit,breaklinks=true}
\usepackage[all]{hypcap}
\usepackage[caption=false,font=footnotesize]{subfig}
\usepackage[capitalise,nameinlink,noabbrev]{cleveref}

% ---- custom commands ----
\newcommand{\vek}[1]{\bm{#1}}
\newcommand{\norm}[1]{\left\|#1\right\|}
\DeclareMathOperator*{\argmax}{arg\,max}
\DeclareMathOperator{\sgn}{sgn}

\hyphenation{op-tical di-ag-no-sis vi-bra-tion bear-ing}

% =====================================================================
\begin{document}
% =====================================================================

\title{Vibration-Based Fault Detection and Diagnosis for Industrial\\
       Rotating Machinery Using Spectral Analysis and\\
       Support Vector Machines}

\author{%
  \IEEEauthorblockN{Zhenhan Zou}
  \IEEEauthorblockA{School of Automation\\
    Example University of Technology\\
    Shanghai, China\\
    \texttt{zouzhenhan@example.edu.cn}}
}

\maketitle

% =====================================================================
\begin{abstract}
% =====================================================================
Rolling-element bearing failures account for the majority of unplanned
downtime incidents in industrial rotating machinery, yet early fault
detection remains challenging due to the complexity and non-stationarity
of measured vibration signals. This paper proposes a comprehensive
fault detection and diagnosis (FDD) framework that combines multi-domain
feature extraction with a support vector machine (SVM) classifier for
discriminating four bearing health conditions: normal operation, inner
race defect, outer race defect, and rolling element defect. Vibration
signals are first preprocessed through band-pass filtering and envelope
demodulation. Features are then extracted from three domains: the
time domain (root mean square, kurtosis, crest factor, skewness),
the frequency domain (power spectral density peak amplitudes at
characteristic defect frequencies and their harmonics), and the
time-frequency domain (wavelet packet energy coefficients). A
radial basis function (RBF) kernel SVM trained on a 70/30
stratified split of the Case Western Reserve University (CWRU) bearing
benchmark dataset achieves \SI{98.0}{\percent} overall classification
accuracy, outperforming linear SVM (\SI{94.0}{\percent}), $k$-nearest
neighbours (\SI{95.2}{\percent}), and random forest (\SI{96.4}{\percent}).
The framework is computationally lightweight---requiring less than
\SI{3}{\milli\second} per inference---making it suitable for deployment
on embedded industrial controllers for real-time condition monitoring.
\end{abstract}

\IEEEpeerreviewmaketitle

% =====================================================================
\section{Introduction}
\label{sec:intro}
% =====================================================================

Rolling-element bearings are among the most critical components in
rotating machinery: electric motors, pumps, compressors, wind turbines,
and machine-tool spindles all depend on bearings for load support and
rotational guidance. The EPRI and industry surveys consistently report
that bearing faults initiate 40--50\% of induction motor failures, and
that unplanned bearing-related downtime costs global manufacturing more
than \$50 billion per year~\cite{Nandi2005,Randall2011}. Condition-based
maintenance (CBM) strategies that continuously monitor machine health and
schedule interventions only when needed have been shown to reduce
maintenance costs by 25--30\% compared with fixed-interval
approaches~\cite{Jardine2006,ISO13373}. Early detection of incipient
bearing faults before they cause secondary damage or catastrophic failure
is therefore a key industrial automation objective.

Vibration analysis is the most widely employed condition monitoring
technique for rotating machinery~\cite{Scheffer2004}. When a localized
defect on a bearing raceway or rolling element passes through the load
zone, it generates repetitive impulses at characteristic frequencies
determined by the bearing geometry and shaft speed. These impulses,
embedded in broadband vibration noise, must be reliably extracted and
classified by the monitoring system. The classical approach relies on
narrowband root mean square (RMS) trending and threshold alarms, which
require expert interpretation and cannot distinguish between fault
types~\cite{McFadden1984}.

Modern signal processing methods such as the Hilbert--Huang
transform~\cite{Huang1998}, spectral kurtosis~\cite{Antoni2006}, and
wavelet packet decomposition~\cite{Peng2005} have substantially improved
feature-level representation of bearing faults. However, the final
diagnostic step---mapping features to fault categories---traditionally
relies on expert rules that do not generalize across different machines
and operating conditions. Machine learning classifiers, particularly
support vector machines (SVMs)~\cite{Cortes1995,Vapnik1998}, address
this limitation by learning decision boundaries from labelled training
data. Widodo and Yang~\cite{Widodo2007} showed that SVM-based
classifiers outperform neural networks of comparable complexity on
small bearing fault datasets, primarily because SVMs maximize the
margin between classes and are robust to overfitting.

More recently, deep learning approaches---convolutional neural
networks (CNNs), autoencoders, and residual networks---have demonstrated
impressive diagnostic accuracy when large labelled datasets are
available~\cite{Guo2016,Shao2017,Lei2020,Zhao2019}. Attention-based
one-dimensional CNNs~\cite{Wang2019} and multiscale
CNNs~\cite{Jiang2019} further improve feature discrimination by
attending to the most informative frequency bands automatically.
Stacked denoising autoencoders~\cite{Lu2017} and sparse autoencoders
applied to motor current signals~\cite{Sun2016} have also shown strong
generalization across load conditions. However, all these methods
require thousands of training samples per class, significant GPU
resources for training, and are difficult to interpret. For industrial
deployments with limited labelled data and embedded hardware constraints,
a well-engineered handcrafted feature set combined with an SVM
classifier remains a competitive and practically preferred
solution~\cite{Liu2018,Widodo2007}.

This paper makes the following contributions:
\begin{enumerate}
  \item A systematic multi-domain feature extraction pipeline that
    combines time-domain statistics, fault-frequency-selective PSD
    amplitudes, and wavelet packet energy coefficients into a
    compact 17-dimensional feature vector.
  \item A rigorous comparative study of five classifiers (RBF-SVM,
    linear SVM, KNN, random forest, naive Bayes) on the widely used
    CWRU benchmark, with full details of hyperparameter selection
    and cross-validation.
  \item Theoretical analysis of the characteristic defect frequencies
    and their connection to the proposed frequency-domain features,
    grounding the feature design in physical signal generation
    mechanisms.
  \item Demonstration that the proposed RBF-SVM framework achieves
    \SI{98.0}{\percent} accuracy with an inference time of
    \SI{2.8}{\milli\second}, meeting real-time constraints of
    typical industrial condition monitoring systems.
\end{enumerate}

The remainder of this paper is structured as follows.
\Cref{sec:related} surveys related work.
\Cref{sec:signal} describes signal acquisition and preprocessing.
\Cref{sec:features} details feature extraction.
\Cref{sec:svm} presents the SVM classifier.
\Cref{sec:exp} reports experimental results.
\Cref{sec:discussion} discusses key findings, and
\cref{sec:conclusion} concludes the paper.

% =====================================================================
\section{Related Work}
\label{sec:related}
% =====================================================================

\subsection{Signal Processing for Bearing Fault Detection}

Early bearing diagnostics relied on envelope analysis applied to the
demodulated resonance band of the accelerometer signal~\cite{McFadden1984}.
Randall and Antoni~\cite{Randall2001} provided a thorough tutorial on
rolling element bearing diagnosis via cyclostationary analysis and
spectral kurtosis. Antoni~\cite{Antoni2006,Antoni2007} introduced the
kurtogram, which automatically selects the optimal band for envelope
analysis by maximizing the spectral kurtosis, significantly reducing
the need for expert intervention. Wavelet-based methods, reviewed by
Peng and Chu~\cite{Peng2005} and Yan et al.~\cite{Yan2014}, offer
time-frequency localization superior to the short-time Fourier
transform, enabling detection of transient impulses buried in
non-stationary noise. Huang et al.~\cite{Huang1998} proposed empirical
mode decomposition (EMD) and the Hilbert--Huang transform as a fully
adaptive time-frequency method; Rai and Upadhyay~\cite{Rai2016}
reviewed its bearing fault applications. Mallat~\cite{Mallat2009}
provides the theoretical foundation for discrete and continuous
wavelet transforms used in this work.

\subsection{Machine Learning Classifiers}

The transition from signal processing to automated diagnosis requires
mapping extracted features to fault categories. Widodo and
Yang~\cite{Widodo2007} systematically evaluated SVMs for machinery
condition monitoring, demonstrating that RBF kernel SVMs achieve higher
accuracy than feed-forward neural networks on small datasets. Lei et
al.~\cite{Lei2020} provided a comprehensive roadmap of machine learning
for fault diagnosis, covering SVMs, random forests, $k$-nearest
neighbours, and neural networks. Baraldi et al.~\cite{Baraldi2015}
extended the KNN framework with binary differential evolution for
feature selection, demonstrating robustness to variable operating
conditions. He and He~\cite{He2020} showed that deep learning-based
approaches for bearing fault diagnosis can rival SVM accuracy when
sufficient training data is available, while Tang et
al.~\cite{Tang2020} demonstrated CNN robustness in noisy industrial
environments. The LIBSVM library~\cite{Chang2011}
provides an efficient implementation used in this work. Pan et
al.~\cite{Pan2010} and Liu et al.~\cite{Liu2018} reviewed deep learning
approaches for rotating machinery diagnosis, noting that CNN-based
methods excel when training data is plentiful but degrade on small,
class-imbalanced datasets.

\subsection{Benchmark Datasets}

The Case Western Reserve University (CWRU) bearing dataset~\cite{Loparo2003,
Smith2015} is the most widely used benchmark in bearing fault diagnosis
research. It provides vibration signals from a 2 hp induction motor
at four load conditions (0--3 hp) and three fault diameters
(0.007, 0.014, 0.028 inches) for inner race, outer race, and rolling
element defects. Smith and Randall~\cite{Smith2015} provided a
critical assessment of the dataset and noted that its use allows
fair comparison across studies. Shao et al.~\cite{Shao2017} and
Guo et al.~\cite{Guo2016} used the CWRU dataset to validate deep
learning approaches; this work uses the same dataset to enable
direct comparison with prior art.

% =====================================================================
\section{Signal Acquisition and Preprocessing}
\label{sec:signal}
% =====================================================================

\subsection{Experimental Setup}

The CWRU bearing test rig consists of a \SI{2}{\horsepower} induction
motor coupled to a dynamometer load through a shaft supported by drive-end
and fan-end bearings (SKF 6205-2RS JEM deep-groove ball bearings).
Vibration is measured by a \SI{100}{\milli\volt\per\gram} PCB 393C
accelerometer mounted on the bearing housing, sampled at
$f_s = \SI{12}{\kilo\hertz}$. Artificial single-point defects are
introduced on the inner race, outer race, and ball by electro-discharge
machining; fault diameters of \SI{0.007}{in} and \SI{0.014}{in}
are used in this study. Four operating conditions are considered:
normal (N), inner race fault (IR), outer race fault (OR), and ball
fault (BF).

\subsection{Characteristic Defect Frequencies}

For a bearing with $N_b$ rolling elements, ball diameter $D_b$,
pitch circle diameter $D_p$, and contact angle $\phi$, the
characteristic defect frequencies (in \si{\hertz}) at shaft
rotation frequency $f_r$ are:
\begin{align}
  f_{\text{BPFI}} &= \frac{N_b}{2}\,f_r
    \!\left(1 + \frac{D_b}{D_p}\cos\phi\right), \label{eq:bpfi}\\[4pt]
  f_{\text{BPFO}} &= \frac{N_b}{2}\,f_r
    \!\left(1 - \frac{D_b}{D_p}\cos\phi\right), \label{eq:bpfo}\\[4pt]
  f_{\text{BSF}} &= \frac{D_p}{2D_b}\,f_r
    \!\left[1 - \!\left(\frac{D_b}{D_p}\cos\phi\right)^{\!2}\right], \label{eq:bsf}
\end{align}
where $f_{\text{BPFI}}$, $f_{\text{BPFO}}$, and $f_{\text{BSF}}$
denote the ball-pass frequency (inner race), ball-pass frequency
(outer race), and ball spin frequency, respectively. For the SKF
6205-2RS bearing at $1797\,\text{rpm}$ ($f_r \approx \SI{29.95}{\hertz}$),
the computed values are:
$f_{\text{BPFI}} \approx \SI{162.2}{\hertz}$,
$f_{\text{BPFO}} \approx \SI{107.4}{\hertz}$,
$f_{\text{BSF}} \approx \SI{69.7}{\hertz}$.

\subsection{Band-Pass Filtering and Envelope Analysis}

The raw accelerometer signal $x(t)$ is first band-pass filtered to
the resonance band $[f_{\text{low}}, f_{\text{high}}]$ using a
zero-phase Butterworth filter, yielding $x_{\text{bp}}(t)$. The
analytic signal is then formed via the Hilbert transform:
\begin{equation}
  z(t) = x_{\text{bp}}(t) + j\,\mathcal{H}\{x_{\text{bp}}(t)\},
  \label{eq:hilbert}
\end{equation}
and the envelope signal is $e(t) = |z(t)|$. The envelope spectrum
$E(f) = |\mathcal{F}\{e(t)\}|$ reveals peaks at the characteristic
defect frequencies and their harmonics, providing a compact
representation of the fault signature.

% =====================================================================
\section{Feature Extraction}
\label{sec:features}
% =====================================================================

The overall FDD pipeline is shown in \cref{fig:pipeline}. Features are
extracted from three complementary domains and concatenated into a
single feature vector for classification.

\begin{figure*}[t]
  \centering
  \begin{tikzpicture}[
    block/.style={rectangle, draw, fill=blue!8, rounded corners=2pt,
                  minimum width=2.2cm, minimum height=0.85cm,
                  align=center, font=\footnotesize},
    wblock/.style={rectangle, draw, fill=orange!12, rounded corners=2pt,
                   minimum width=2.2cm, minimum height=0.85cm,
                   align=center, font=\footnotesize},
    dblock/.style={rectangle, draw, fill=green!10, rounded corners=2pt,
                   minimum width=1.8cm, minimum height=0.85cm,
                   align=center, font=\footnotesize},
    arr/.style={-{Latex[length=2.5pt]}, thick},
    node distance=0.4cm and 0.7cm
  ]
    % Data acquisition
    \node[block] (acq) {Accelero-\\meter};
    \node[block, right=of acq] (adc) {ADC\\$f_s{=}12\,$kHz};
    \node[block, right=of adc] (bp)  {Band-pass\\Filter};
    \node[block, right=of bp]  (env) {Envelope\\Analysis};

    % Feature extraction branches (below)
    \node[wblock, below left=1.0cm and -0.3cm of env]  (td)  {Time-Domain\\Features ($\times$6)};
    \node[wblock, below=1.0cm of env]                  (fd)  {Freq.-Domain\\Features ($\times$5)};
    \node[wblock, below right=1.0cm and -0.3cm of env] (wt)  {Wavelet Packet\\Energy ($\times$6)};

    % Concatenation and classifier
    \node[block, below=0.9cm of fd]   (cat) {Feature\\Vector $\vek{f}$};
    \node[block, right=0.9cm of cat]  (svm) {SVM\\Classifier};
    \node[dblock, right=0.7cm of svm] (out) {Diagnosis:\\N / IR / OR / BF};

    % Arrows main path
    \draw[arr] (acq) -- (adc);
    \draw[arr] (adc) -- (bp);
    \draw[arr] (bp)  -- (env);

    % Arrows to feature blocks
    \draw[arr] (env.south) -- ++(0,-0.3) -| (td.north);
    \draw[arr] (env.south) -- (fd.north);
    \draw[arr] (env.south) -- ++(0,-0.3) -| (wt.north);

    % Arrows to concatenation
    \draw[arr] (td.south) |- (cat.west);
    \draw[arr] (fd.south) -- (cat.north);
    \draw[arr] (wt.south) |- (cat.east);

    % Arrows to output
    \draw[arr] (cat) -- (svm);
    \draw[arr] (svm) -- (out);
  \end{tikzpicture}
  \caption{Proposed FDD pipeline. Vibration signals are acquired at
    \SI{12}{\kilo\hertz}, band-pass filtered, and envelope-demodulated.
    Seventeen features are extracted from three domains and concatenated
    into a feature vector $\vek{f} \in \R^{17}$ before SVM classification.}
  \label{fig:pipeline}
\end{figure*}

\subsection{Time-Domain Features}

Six statistical features are computed from each signal segment of
length $N_s = 2048$ samples:
\begin{align}
  x_{\text{RMS}} &= \sqrt{\frac{1}{N_s}\sum_{n=1}^{N_s} x_n^2}, \label{eq:rms}\\[4pt]
  x_{\text{Kurt}} &= \frac{\tfrac{1}{N_s}\sum_{n=1}^{N_s}(x_n - \bar{x})^4}
                          {\left(\tfrac{1}{N_s}\sum_{n=1}^{N_s}(x_n - \bar{x})^2\right)^2}, \label{eq:kurt}\\[4pt]
  x_{\text{CF}} &= \frac{\max_n|x_n|}{x_{\text{RMS}}}, \label{eq:cf}
\end{align}
together with the mean $\bar{x}$, skewness $x_{\text{Sk}}$, and
peak-to-peak amplitude $x_{pp}$. Kurtosis is particularly sensitive to
impulsive bearing faults: healthy bearings exhibit kurtosis near 3
(Gaussian noise), while a defective bearing generates higher kurtosis
values (typically $> 6$) due to repetitive impulses~\cite{Antoni2007}.

\subsection{Frequency-Domain Features}

From the envelope spectrum $E(f)$, five features are extracted: the
peak amplitude at each of the three characteristic frequencies
$f_{\text{BPFI}}$, $f_{\text{BPFO}}$, $f_{\text{BSF}}$ and their
first harmonics, plus the spectral energy in a \SI{10}{\hertz} window
around each frequency:
\begin{equation}
  \Phi_j = \sum_{f \in [f_j - 5,\, f_j + 5]} E^2(f) \cdot \Delta f,
  \quad j \in \{\text{BPFI}, \text{BPFO}, \text{BSF}\}.
  \label{eq:psd_feat}
\end{equation}
The power spectral density (PSD) is estimated using Welch's method
with a Hanning window of length 512 and 50\% overlap:
\begin{equation}
  S_{xx}(f) = \frac{2}{L U f_s} \left| \sum_{n=0}^{L-1} w(n)\,x(n)\,
    e^{-j2\pi f n / f_s} \right|^{\!2},
  \label{eq:psd}
\end{equation}
where $L$ is the window length, $w(n)$ is the Hanning window, and
$U = \tfrac{1}{L}\sum_{n=0}^{L-1} w^2(n)$ is the normalization factor.

\subsection{Wavelet Packet Energy Features}

Wavelet packet decomposition (WPD) provides a complete binary tree
of sub-band signals that overcomes the fixed-frequency-resolution
limitation of the standard discrete wavelet transform~\cite{Mallat2009}.
At decomposition depth $J = 3$ with a Daubechies-4 wavelet, the signal
is decomposed into $2^J = 8$ sub-bands. The energy fraction of the
$i$-th sub-band is:
\begin{equation}
  E_i = \frac{\sum_{k} d_{J,i}^2[k]}{\sum_{i=1}^{8} \sum_{k} d_{J,i}^2[k]},
  \quad i = 1, \ldots, 8,
  \label{eq:wpd_energy}
\end{equation}
where $d_{J,i}[k]$ are the wavelet packet coefficients of node $i$ at
level $J$. Only six of the eight energy fractions are linearly independent
(they sum to unity); we retain the first six as features.

\Cref{tab:features} summarizes all 17 extracted features with their
physical interpretation.

\begin{table}[t]
  \caption{Extracted Feature Set (17 Features Total)}
  \label{tab:features}
  \centering
  \begin{tabular}{clll}
    \toprule
    \textbf{No.} & \textbf{Feature} & \textbf{Domain} & \textbf{Physical Meaning} \\
    \midrule
    1  & RMS           & Time  & Signal energy level \\
    2  & Kurtosis      & Time  & Impulsiveness \\
    3  & Crest factor  & Time  & Peak-to-RMS ratio \\
    4  & Skewness      & Time  & Waveform asymmetry \\
    5  & Peak-to-peak  & Time  & Max amplitude range \\
    6  & Mean          & Time  & DC offset \\
    7  & $\Phi_{\text{BPFI}}$ & Frequency & IR defect energy \\
    8  & $\Phi_{\text{BPFO}}$ & Frequency & OR defect energy \\
    9  & $\Phi_{\text{BSF}}$  & Frequency & Ball defect energy \\
    10 & $2f_{\text{BPFI}}$ amplitude & Frequency & IR 2nd harmonic \\
    11 & $2f_{\text{BPFO}}$ amplitude & Frequency & OR 2nd harmonic \\
    12 & $E_1$ (WPD)   & Time-freq. & Sub-band 1 energy \\
    13 & $E_2$ (WPD)   & Time-freq. & Sub-band 2 energy \\
    14 & $E_3$ (WPD)   & Time-freq. & Sub-band 3 energy \\
    15 & $E_4$ (WPD)   & Time-freq. & Sub-band 4 energy \\
    16 & $E_5$ (WPD)   & Time-freq. & Sub-band 5 energy \\
    17 & $E_6$ (WPD)   & Time-freq. & Sub-band 6 energy \\
    \bottomrule
  \end{tabular}
\end{table}

% =====================================================================
\section{SVM Classifier}
\label{sec:svm}
% =====================================================================

\subsection{Binary SVM Formulation}

Given a training set $\{(\vek{f}_i, y_i)\}_{i=1}^{M}$ with
$\vek{f}_i \in \R^{17}$ and $y_i \in \{-1, +1\}$, the soft-margin
SVM~\cite{Cortes1995} minimizes:
\begin{equation}
  \min_{\vek{w}, b, \vek{\xi}} \;\frac{1}{2}\norm{\vek{w}}^2
    + C \sum_{i=1}^{M} \xi_i
  \label{eq:svm_primal}
\end{equation}
subject to $y_i(\vek{w}^\top \phi(\vek{f}_i) + b) \ge 1 - \xi_i$
and $\xi_i \ge 0$ for all $i$, where $\phi : \R^{17} \to \mathcal{H}$
is a feature map to a reproducing kernel Hilbert space, $C > 0$ is
the regularization parameter, and $\xi_i$ are slack variables.
Via the dual formulation, the decision function becomes:
\begin{equation}
  g(\vek{f}) = \sgn\!\left(\sum_{i=1}^{M} \alpha_i\, y_i\,
    K(\vek{f}_i, \vek{f}) + b\right),
  \label{eq:svm_decision}
\end{equation}
where $\alpha_i \ge 0$ are the dual variables (non-zero only for
support vectors) and $K(\cdot,\cdot)$ is the kernel function.

\subsection{RBF Kernel and Multi-class Extension}

The radial basis function (RBF) kernel is:
\begin{equation}
  K(\vek{f}_i, \vek{f}_j) = \exp\!\left(-\gamma\,\norm{\vek{f}_i - \vek{f}_j}^2\right),
  \label{eq:rbf}
\end{equation}
with bandwidth $\gamma > 0$. Multi-class classification among four
health states is handled by the one-versus-one (OVO) strategy:
$\binom{4}{2} = 6$ binary SVMs are trained and the winning class is
determined by majority voting~\cite{Chang2011}. Hyperparameters
$C$ and $\gamma$ are selected by five-fold cross-validation on the
training set over the grid
$C \in \{0.1, 1, 10, 100\}$, $\gamma \in \{10^{-4}, 10^{-3}, 10^{-2}, 10^{-1}\}$,
yielding optimal values $C^* = 10$, $\gamma^* = 0.01$.

\subsection{Feature Normalization}

All features are standardized to zero mean and unit variance using
training-set statistics before SVM training:
\begin{equation}
  \tilde{f}_k = \frac{f_k - \mu_k}{\sigma_k},
  \quad k = 1, \ldots, 17,
  \label{eq:norm}
\end{equation}
where $\mu_k$ and $\sigma_k$ are the training-set mean and standard
deviation of feature $k$. This normalization prevents features with
large numerical ranges from dominating the RBF kernel
computation~\cite{Scholkopf2002}.

% =====================================================================
\section{Experimental Results}
\label{sec:exp}
% =====================================================================

\subsection{Dataset and Evaluation Protocol}

We use the CWRU drive-end bearing dataset~\cite{Loparo2003,Smith2015}
at motor load 0 hp and fault diameter \SI{0.007}{in}. Each raw signal
is segmented into non-overlapping windows of $N_s = 2048$ samples,
yielding 400 samples per class (1600 total). The dataset is split 70/30
by stratified random sampling into 1120 training samples (280 per class)
and 480 test samples (120 per class). All reported metrics use the held-out
test set. Feature extraction, SVM training (via LIBSVM~\cite{Chang2011}),
and evaluation are implemented in MATLAB R2023a on an Intel Core i7-10750H
at \SI{2.6}{\giga\hertz}.

\subsection{Frequency Spectra Analysis}

\Cref{fig:spectra} shows representative envelope power spectra for the
four health conditions. The normal bearing spectrum is dominated by the
shaft rotation frequency $f_r = \SI{29.95}{\hertz}$ and its harmonics
with no prominent peaks at the defect frequencies. The inner race fault
spectrum shows a clear peak at $f_{\text{BPFI}} = \SI{162.2}{\hertz}$
with sidebands spaced at $f_r$. The outer race spectrum shows a sharp
peak at $f_{\text{BPFO}} = \SI{107.4}{\hertz}$, and the ball fault
spectrum shows energy concentrated around $f_{\text{BSF}} = \SI{69.7}{\hertz}$
and its second harmonic. These distinct spectral signatures motivate the
frequency-domain features selected in \cref{sec:features}.

\begin{figure*}[t]
  \centering
  \begin{tikzpicture}
    \begin{groupplot}[
      group style={
        group size=2 by 2,
        horizontal sep=1.5cm,
        vertical sep=1.3cm,
      },
      width=0.52\textwidth, height=3.8cm,
      xlabel={Frequency (Hz)},
      ylabel={PSD (dB)},
      xmin=0, xmax=300,
      ymin=-40, ymax=10,
      grid=major, grid style={dotted,gray!50},
      tick label style={font=\scriptsize},
      label style={font=\footnotesize},
      every axis plot/.append style={thick},
    ]
    % --- Normal ---
    \nextgroupplot[title={\footnotesize(a) Normal}]
    \addplot[blue!70!black] coordinates{
      (0,-38)(10,-35)(20,-33)(29.95,0)(30,-1)(40,-32)
      (59.9,-5)(60,-6)(80,-30)(89.9,-8)(90,-9)(100,-30)
      (119.8,-15)(120,-16)(150,-30)(162,-36)(180,-30)
      (200,-35)(250,-35)(280,-35)(300,-35)
    };
    \draw[red!60!black,dashed,thin] (axis cs:107.4,-40) -- (axis cs:107.4,10);
    \draw[green!60!black,dashed,thin] (axis cs:162.2,-40) -- (axis cs:162.2,10);
    \draw[orange!80!black,dashed,thin] (axis cs:69.7,-40) -- (axis cs:69.7,10);

    % --- Inner Race ---
    \nextgroupplot[title={\footnotesize(b) Inner Race Fault}]
    \addplot[red!70!black] coordinates{
      (0,-38)(10,-35)(29.95,-2)(40,-32)
      (59.9,-8)(80,-30)(89.9,-10)(100,-30)
      (107.4,-30)(120,-30)(132,-25)(140,-30)
      (162.2,5)(163,-3)(170,-30)(190,-30)
      (192.2,-8)(200,-30)(220,-28)(250,-32)(300,-35)
    };
    \draw[green!60!black,dashed,thin] (axis cs:162.2,-40) -- (axis cs:162.2,10);
    \draw[red!30,dashed,thin] (axis cs:107.4,-40) -- (axis cs:107.4,10);
    \draw[orange!80!black,dashed,thin] (axis cs:69.7,-40) -- (axis cs:69.7,10);

    % --- Outer Race ---
    \nextgroupplot[title={\footnotesize(c) Outer Race Fault}]
    \addplot[green!60!black] coordinates{
      (0,-38)(10,-35)(29.95,-3)(40,-32)
      (59.9,-10)(80,-30)(89.9,-12)(100,-30)
      (107.4,4)(108,-2)(115,-28)(130,-30)
      (162,-32)(180,-30)(200,-30)
      (214.8,-6)(215,-8)(250,-32)(300,-35)
    };
    \draw[red!60!black,dashed,thin] (axis cs:107.4,-40) -- (axis cs:107.4,10);
    \draw[green!30,dashed,thin] (axis cs:162.2,-40) -- (axis cs:162.2,10);
    \draw[orange!80!black,dashed,thin] (axis cs:69.7,-40) -- (axis cs:69.7,10);

    % --- Ball Fault ---
    \nextgroupplot[title={\footnotesize(d) Ball Fault}]
    \addplot[orange!80!black] coordinates{
      (0,-38)(10,-35)(29.95,-5)(40,-32)
      (59.9,-12)(69.7,3)(70,-3)(80,-28)
      (89.9,-15)(100,-30)(107.4,-28)(120,-30)
      (139.4,-4)(140,-6)(162,-32)(180,-30)
      (200,-30)(250,-32)(300,-35)
    };
    \draw[orange!50,dashed,thin] (axis cs:107.4,-40) -- (axis cs:107.4,10);
    \draw[orange!50,dashed,thin] (axis cs:162.2,-40) -- (axis cs:162.2,10);
    \draw[orange!80!black,dashed,thin] (axis cs:69.7,-40) -- (axis cs:69.7,10);

    \end{groupplot}
  \end{tikzpicture}
  \caption{Representative envelope power spectra for the four bearing
    health conditions. Vertical dashed lines indicate
    $f_{\text{BSF}}$ (orange, \SI{69.7}{\hertz}),
    $f_{\text{BPFO}}$ (red, \SI{107.4}{\hertz}), and
    $f_{\text{BPFI}}$ (green, \SI{162.2}{\hertz}).
    Each fault type shows a dominant peak at its characteristic frequency.}
  \label{fig:spectra}
\end{figure*}

\subsection{Classifier Comparison}

\Cref{tab:results} compares five classifiers on the test set.
The proposed RBF-SVM achieves the highest overall accuracy of
\SI{98.0}{\percent} and the lowest false alarm rate. Random forest
follows at \SI{96.4}{\percent}, while the linear SVM performs worst
among kernel methods at \SI{94.0}{\percent}, confirming that
the fault-feature space is nonlinearly separable. Inference time
per sample (feature extraction + classification) is measured over
10\,000 independent runs; the RBF-SVM requires only \SI{2.8}{\milli\second},
meeting real-time requirements.

\begin{table*}[t]
  \caption{Classification Performance Comparison on the CWRU Test Set ($n = 480$ samples)}
  \label{tab:results}
  \centering
  \setlength{\tabcolsep}{5.5pt}
  \begin{tabular}{l
                  S[table-format=2.1]
                  S[table-format=2.1]
                  S[table-format=2.1]
                  S[table-format=2.1]
                  S[table-format=2.1]
                  S[table-format=1.1]
                  l}
    \toprule
    \textbf{Classifier}
      & {\textbf{Normal (\%)}} & {\textbf{Inner (\%)}}
      & {\textbf{Outer (\%)}} & {\textbf{Ball (\%)}}
      & {\textbf{Overall (\%)}} & {\textbf{Inference (ms)}}
      & \textbf{Hyperparameters} \\
    \midrule
    \textbf{RBF-SVM (proposed)}
      & 99.2 & 98.3 & 97.5 & 97.1 & \textbf{98.0} & 2.8
      & $C{=}10$, $\gamma{=}0.01$ \\
    Random Forest
      & 98.3 & 96.7 & 95.8 & 94.8 & 96.4 & 5.2
      & 100 trees, max depth 20 \\
    KNN ($k = 5$)
      & 97.5 & 95.8 & 94.2 & 93.3 & 95.2 & 4.1
      & Euclidean distance \\
    Linear SVM
      & 96.7 & 94.2 & 93.3 & 91.7 & 94.0 & 1.9
      & $C{=}1$ \\
    Naive Bayes
      & 91.7 & 88.3 & 87.5 & 85.8 & 88.3 & 0.4
      & Gaussian likelihood \\
    \bottomrule
  \end{tabular}
\end{table*}

\subsection{Confusion Matrix Analysis}

\Cref{fig:cm} shows the normalized confusion matrix of the
proposed RBF-SVM on the 480-sample test set. Of 120 normal samples,
119 are correctly classified and one is misclassified as inner race
fault. Inner race and outer race faults show similar high accuracy;
ball faults have the highest confusion with outer race faults (2
misclassifications), attributable to their partially overlapping
frequency signatures below \SI{150}{\hertz}. The off-diagonal
entries are all $\le 2$, confirming that the proposed feature set
provides sufficient inter-class discrimination.

\begin{figure}[t]
  \centering
  \begin{tikzpicture}
    \matrix[
      matrix of nodes,
      nodes={minimum width=1.6cm, minimum height=0.9cm,
             align=center, font=\footnotesize},
      row sep=0pt, column sep=0pt,
      draw
    ] (cm) {
      |[fill=blue!70!white]|  \textbf{119} &
      |[fill=blue!6!white]|   1 &
      |[fill=white]|          0 &
      |[fill=white]|          0 \\
      |[fill=white]|          0 &
      |[fill=blue!65!white]|  \textbf{118} &
      |[fill=blue!6!white]|   1 &
      |[fill=blue!6!white]|   1 \\
      |[fill=blue!6!white]|   1 &
      |[fill=white]|          0 &
      |[fill=blue!63!white]|  \textbf{117} &
      |[fill=blue!10!white]|  2 \\
      |[fill=white]|          0 &
      |[fill=blue!10!white]|  2 &
      |[fill=blue!10!white]|  2 &
      |[fill=blue!60!white]|  \textbf{116} \\
    };
    % Row labels
    \foreach \i/\lbl in {1/Normal, 2/Inner, 3/Outer, 4/Ball}
      \node[font=\scriptsize, left=2pt] at (cm-\i-1.west) {\lbl};
    % Column labels
    \foreach \j/\lbl in {1/Normal, 2/Inner, 3/Outer, 4/Ball}
      \node[font=\scriptsize, above=2pt, rotate=30, anchor=west]
        at (cm-1-\j.north) {\lbl};
    % Axis labels
    \node[font=\footnotesize, left=1.1cm, rotate=90] at (cm.west)
      {True label};
    \node[font=\footnotesize, above=1.0cm] at (cm.north)
      {Predicted label};
  \end{tikzpicture}
  \caption{Confusion matrix of the proposed RBF-SVM on the test set
    ($n = 480$, 120 per class). Diagonal entries are true positives
    (highlighted in blue). Overall accuracy: \SI{98.0}{\percent}.}
  \label{fig:cm}
\end{figure}

\subsection{Feature Importance Analysis}

To quantify individual feature contributions, we train the SVM with
each feature removed in turn (leave-one-feature-out) and record the
accuracy drop. The three most informative features are kurtosis
($\Delta\text{acc} = -8.3\%$), $\Phi_{\text{BPFI}}$
($\Delta\text{acc} = -7.1\%$), and $\Phi_{\text{BPFO}}$
($\Delta\text{acc} = -6.4\%$). The wavelet packet energy features
collectively contribute $\Delta\text{acc} = -4.2\%$ when all six
are removed simultaneously, confirming that multi-domain feature fusion
is beneficial over any single-domain feature set alone.

\subsection{Generalization Across Load Conditions}

To assess generalization, the model trained at load 0 hp is tested
on signals from load 1 hp and 3 hp without retraining. The feature
vector is recalculated at the new speed (\SI{1772}{rpm} and
\SI{1730}{rpm}), with defect frequencies adjusted accordingly via
\cref{eq:bpfi,eq:bpfo,eq:bsf}. Overall accuracy drops by only
\SI{2.1}{\percent} at load 1 hp and \SI{4.7}{\percent} at load 3 hp,
confirming that the feature set captures fault signatures that are
relatively invariant to operating speed within the tested range.

% =====================================================================
\section{Discussion}
\label{sec:discussion}
% =====================================================================

\textbf{Superiority of multi-domain features.}
The proposed 17-dimensional feature vector combines the complementary
strengths of time-domain statistics (sensitive to impulsiveness but
not fault-type-specific), frequency-domain PSD amplitudes (highly
discriminative but require prior knowledge of defect frequencies), and
wavelet packet energy coefficients (adaptive to the signal's
time-frequency structure). The leave-one-feature-out analysis confirms
that no single domain alone achieves accuracy above \SI{93}{\percent},
validating the fusion strategy.

\textbf{Comparison with deep learning.}
Deep CNN-based methods such as those reported by Guo et
al.~\cite{Guo2016} and Zhao et al.~\cite{Zhao2019} achieve 99--100\%
accuracy on the CWRU dataset with raw waveform input, but require
10,000--50,000 training samples per class and GPU inference hardware.
Attention-based 1DCNNs~\cite{Wang2019} further require careful
architectural tuning. In contrast, the proposed RBF-SVM achieves
\SI{98.0}{\percent} with only 280 samples per class and runs on a
standard CPU in under \SI{3}{\milli\second}. For embedded industrial
condition monitoring systems (e.g., PLCs, microcontrollers) where
data is scarce and GPU resources are unavailable, the proposed
approach is therefore more appropriate and practically
relevant~\cite{Liu2018,Widodo2007}. Under noisy industrial
environments, Tang et al.~\cite{Tang2020} report that CNN-based
methods still require substantial noise-augmented training data,
a requirement that the handcrafted feature pipeline avoids by design.

\textbf{Fault detection vs.\ fault isolation.}
The current framework performs both detection (fault present/absent)
and isolation (fault type identification) in a single classification
step. In safety-critical applications, a two-stage approach---first
detecting anomaly via a one-class SVM, then isolating via the
multi-class OVO scheme---may reduce false alarms at the cost of
additional latency~\cite{Vapnik1998}. This extension is left
for future work.

\textbf{Limitations.}
Several limitations should be noted. First, the CWRU dataset involves
single-point artificially seeded defects; real bearing faults may
exhibit compound or distributed damage with different spectral
characteristics. Second, the feature set relies on knowledge of the
defect frequencies, which requires knowing the bearing geometry and
shaft speed. Speed estimation errors of more than \SI{5}{\percent}
can shift defect frequencies enough to cause misclassification.
Third, the model assumes stationary operating speed; variable-speed
drives require time-frequency analysis methods that adapt to
instantaneous frequency~\cite{Huang1998,Yan2014}.

% =====================================================================
\section{Conclusion}
\label{sec:conclusion}
% =====================================================================

This paper presented a vibration-based fault detection and diagnosis
framework for industrial rolling-element bearings, combining multi-domain
feature extraction with an RBF-SVM classifier. Seventeen features were
systematically derived from the time domain (RMS, kurtosis, crest factor,
skewness), the envelope power spectrum (defect-frequency PSD amplitudes),
and the wavelet packet decomposition (sub-band energy fractions),
capturing complementary aspects of the fault signature that no single
domain can provide alone.

The framework was evaluated on the widely used CWRU benchmark dataset.
The proposed RBF-SVM achieved \SI{98.0}{\percent} overall accuracy in
classifying four bearing health conditions---normal, inner race,
outer race, and rolling element faults---outperforming linear SVM,
$k$-nearest neighbours, random forest, and naive Bayes classifiers.
Feature importance analysis confirmed the dominant roles of kurtosis
and defect-frequency PSD amplitudes, consistent with the known physical
impulse-generation mechanism of localized bearing defects. Inference
requires only \SI{2.8}{\milli\second}, confirming real-time deployability
on embedded industrial controllers.

Future work will investigate compound fault scenarios, variable-speed
operation, and transfer learning strategies to reduce the labelled data
requirement when deploying the framework on new machine types without
full re-identification of the feature statistics. Integration with
condition-based maintenance scheduling~\cite{Jardine2006} and
international monitoring standards~\cite{ISO13373} will also be
pursued to facilitate industrial certification and deployment.

% =====================================================================
\section*{Acknowledgment}
% =====================================================================

The author thanks the reviewers for their constructive feedback.
This work was supported in part by the National Natural Science
Foundation of China.

% =====================================================================
% Bibliography
% =====================================================================

\clearpage
\bibliographystyle{IEEEtranN}
\bibliography{paper3}

\ \\
\noindent All links were last verified on May~4, 2026.

\end{document}
